%======================================================================
\section{Two steps up the ladder}\label{sec:ladder}
%======================================================================

Before the main construction we climb two steps that show what the criterion needs. The first keeps transport by the powers of one element
and makes it cheap: in an elementary amenable group the far commutators have at most quadratic area, which gives the exponent $1/3$.
Transport by the powers of one element gives exponent $1/(p+1)$ when the far commutators have area of order $k^p$, and along a ray of $H_3$ the
area is at least quadratic~\cite[Proposition~3.4]{DouglasMghirbiA}, so that transport stops at $1/3$. The second changes the kind of transport: Thompson's group $F$ has orbits that grow
exponentially, so words of length $O(\log N)$ reach $N$ sites, and its quadratic Dehn function fills all pairwise commutators cheaply. This comes
within a factor $\log^2r$ in the exponent of the ceiling of Theorem~\ref{thm:criterion}, but in a group that contains $F$, whose soficity is unknown.
Section~\ref{sec:group} keeps the exponential transport and removes the condition.

\subsection{An amenable group at exponent \texorpdfstring{$1/3$}{1/3}}\label{sec:dyadic}

In this subsection and the next, functions are composed from right to left, and a group acting on a set $X$ acts on
$\bigoplus_{x\in X}S_3$ by carrying the copy of $S_3$ at $x$ to the copy at $g(x)$ by the identity of $S_3$.

Let $X=\mathbb Z[\frac12]$, and let $t$ and $u$ be the permutations of $X$ given by $t(x)=x+1$ for $x\in\mathbb Z$, $t(x)=x$ for
$x\notin\mathbb Z$, and $u(x)=2x$. Let $\mathcal B=\langle t,u\rangle$ and $\mathcal G=\bigl(\bigoplus_{x\in X}S_3\bigr)\rtimes\mathcal B$, and let $a,b$ be
generating involutions of the copy of $S_3$ at $0$. Put $s=utu^{-1}$ and $v=tst^{-1}$. Then $s$ adds $2$ to the even integers and $v$ adds $2$
to the odd integers, each fixing every other point, so $t^2=sv$ and $[s,v]=1$. Moreover $u$ and $v$ fix $0$, and $s$ fixes $1=t(0)$. Since $\mathcal B$ acts transitively on $X$, the group $\mathcal G$ is generated by $a$,
$b$, $t$ and $u$. So the
following seventeen words in $a,b,t,u,s,v$, which make up the set $R_{\mathcal G}$, are relations of $\mathcal G$:
\[
a^2,\ b^2,\ (ab)^3,\quad s^{-1}utu^{-1},\ v^{-1}tst^{-1},\ t^2v^{-1}s^{-1},\ [s,v],\quad [u,x],\ [v,x],\ [s,txt^{-1}],\ [x,tyt^{-1}],
\]
where $x,y\in\{a,b\}$.

\begin{theorem}\label{thm:dyadic}
The group $\mathcal G$ is elementary amenable and not solvable, and every commutator $[x,t^kyt^{-k}]$ with $x,y\in\{a,b\}$ and
$k\ge1$ has area at most $5k^2$ with respect to $R_{\mathcal G}$. Consequently the chunk $E_{\mathcal G}$ that Theorem~\ref{thm:criterion}
associates with $R_{\mathcal G}$ has finite sofic profile with
\[
\log\sigma_{E_{\mathcal G}}(r)\ \ge\ \frac1{16}\Big\lfloor\Big(\frac r{75\cdot10^8}\Big)^{1/3}\Big\rfloor\qquad(r>1),
\]
every unitary assignment $\phi$ as in Theorem~\ref{thm:criterion} has $\log D\ge\frac{\Delta^2}{12}\bigl((5K_\Delta e)^{-2/5}-1\bigr)$, and
$\mathcal G$ is not LEF.
\end{theorem}

\begin{proof}
\emph{Amenability.} For $j\in\mathbb Z$ let $t_j=u^jtu^{-j}$; it adds $2^j$ to the points of $2^j\mathbb Z$ and fixes the others. The group $\mathcal K$
generated by all $t_j$ is normalized by $u$ and contains $t$, so it is normal in $\mathcal B$ and $\mathcal B/\mathcal K$ is cyclic. For $n\ge1$ the elements $t_j$ with
$|j|\le n$ move only points of $2^{-n}\mathbb Z$, and they commute with the translation $x\mapsto x+2^n$ of $2^{-n}\mathbb Z$, which has $4^n$
orbits. A permutation of $2^{-n}\mathbb Z$ that commutes with this translation maps each orbit onto an orbit by a translation, so these $t_j$ generate a subgroup of
$\mathbb Z^{4^n}\rtimes\Sym_{4^n}$, which is virtually abelian. Hence $\mathcal K$ is a directed union of virtually abelian groups, $\mathcal B$ is elementary
amenable, and so is $\mathcal G$, because $\bigoplus_XS_3$ is locally finite.

\emph{Not solvable.} For $n\ge1$ the elements $t_0,\dots,t_{n-1}$ preserve $\mathbb Z$ and commute with $x\mapsto x+2^n$, so they act on
$\mathbb Z/2^n$, where they permute the classes modulo $2^i$ for every $i\le n$. So they generate a subgroup $\Gamma_n$ of the iterated wreath
product $\mathsf W_n$ of $n$ copies of $\mathbb Z/2$, which has order $2^{2^n-1}$. On the even classes, identified with $\mathbb Z/2^{n-1}$ by
halving, $t_1,\dots,t_{n-1}$ act as $t_0,\dots,t_{n-2}$ act on $\mathbb Z/2^{n-1}$, and they fix the odd classes; conjugation by $t_0$ carries
this action to the odd classes, and $t_0$ exchanges the two halves. So $|\Gamma_n|\ge2|\Gamma_{n-1}|^2$, and $\Gamma_n=\mathsf W_n$ by induction from $\Gamma_1=\mathsf W_1=\mathbb Z/2$. The derived
subgroup of $\mathsf W_n$ lies in $\mathsf W_{n-1}\times\mathsf W_{n-1}$ and projects onto each factor, since the commutator of $(g,1)$ with the
exchange is $(g^{-1},g)$; by induction the $(n-1)$st derived subgroup of $\mathsf W_n$ is nontrivial. Since $\Gamma_n$ is a quotient of the subgroup
$\langle t_0,\dots,t_{n-1}\rangle$ of $\mathcal B$, the derived length of $\mathcal B$ is at least $n$ for every $n$, so neither $\mathcal B$ nor
$\mathcal G$ is solvable.

\emph{Area.} Let $A_{\mathcal G}(k)$ be the largest area of $[x,t^kyt^{-k}]$ over $x,y\in\{a,b\}$, so that $A_{\mathcal G}(1)=1$. Each of the following moves costs one
cell: replacing $t^2$ by $sv$; exchanging adjacent letters $s^{\pm1}$ and $v^{\pm1}$; replacing $s$ by $utu^{-1}$; moving a letter $x^{\pm1}$,
$x\in\{a,b\}$, across $u^{\pm1}$ or $v^{\pm1}$; and moving $txt^{-1}$ or its inverse across $s^{\pm1}$. The commutator contains four powers
$t^{\pm k}$.

For $k=2m$, write each $t^{\pm2m}$ as $(s^mv^m)^{\pm1}$, with $m$ replacements and $m(m-1)/2$ exchanges each, $2m^2+2m$ cells in all. Moving
$y^{\pm1}$ across $v^{\pm m}$ ($2m$ cells) leaves $[x,s^mys^{-m}]$. Replacing its $4m$ letters $s^{\pm1}$ gives
$[x,ut^mu^{-1}yut^{-m}u^{-1}]$, and moving $y^{\pm1}$ and $x^{\pm1}$ across $u$ (four cells) gives $u[x,t^myt^{-m}]u^{-1}$. Hence
$A_{\mathcal G}(2m)\le A_{\mathcal G}(m)+2m^2+8m+4$.

For $k=2m+1$, put $z=tyt^{-1}$, so that the commutator is freely $[x,t^{2m}zt^{-2m}]$. Write each $t^{\pm2m}$ as $(v^ms^m)^{\pm1}$, now with
$m(m+1)/2$ exchanges each, $2m^2+6m$ cells in all. Moving $z^{\pm1}$ across $s^{\pm m}$ ($2m$ cells) leaves $[x,v^mzv^{-m}]$, moving
$x^{\pm1}$ across $v^{\pm m}$ ($2m$ cells) leaves $v^m[x,z]v^{-m}$, and $[x,z]$ is a relation. Hence $A_{\mathcal G}(2m+1)\le2m^2+10m+1$.

By induction $A_{\mathcal G}(k)\le5k^2$: $A_{\mathcal G}(2)\le15$, $A_{\mathcal G}(m)+2m^2+8m+4\le7m^2+8m+4\le20m^2$ for $m\ge1$, and $2m^2+10m+1\le5(2m+1)^2$.

\emph{Profile.} The words of $R_{\mathcal G}$ have length at most $8$, so Corollary~\ref{cor:cyclic} with $C=5$, $p=2$ and $l=8$ gives the
bounds. Since $\mathcal G$ is amenable, it is sofic and every chunk has finite profile. The profile of $E_{\mathcal G}$ is unbounded, so
$E_{\mathcal G}$ does not embed in a finite group (\cite[Lemma~2.1]{DouglasMghirbiA}), and $\mathcal G$ is not LEF.
\end{proof}

The fillings in the proof have been generated and checked by free reduction for $k\le64$; the largest ratio of cell count to $k^2$ among
them is $15/4$, at $k=2$. In $H_3$ the doubling endomorphism multiplies a reweighted area by $107$ at each step, which gives the exponent $\log107$; in
$\mathcal G$ the dilation $u$ relates distance $2m$ to distance $m$ exactly, at a cost of $O(m^2)$ cells for the two commuting shifts. The orbit of $0$ under $\mathcal B$
grows exponentially, as the orbits of Thompson's group $F$ do, but we do not know whether $\mathcal G$ supports the transport of many
copies used for $F$ below (Section~\ref{sec:open}).

\subsection{Thompson's group \texorpdfstring{$V$}{V}}\label{sec:thompson}

Genevois proved that every Houghton group embeds in Thompson's group $V$~\cite{Genevois2019}. Since $H_3$ is amenable, $V$ therefore has a
chunk whose profile is finite at every $r$ and superpolynomial by~\cite[Theorem~1.1]{DouglasMghirbiA}. A different subgroup of $V$ gives a much larger bound.

Let $\mathbb D$ be the set of dyadic rationals in $(0,1)$, and let Thompson's group $F$ act on $[0,1]$ by its piecewise-linear
homeomorphisms. Let $\mathcal L=\bigl(\bigoplus_{x\in\mathbb D}S_3\bigr)\rtimes F$.

\begin{lemma}[Gap coding]\label{lem:gap-coding}
$\mathcal L$ embeds in $V$.
\end{lemma}

\begin{proof}
Let $\mathsf c$ be the substitution $0\mapsto0$, $1\mapsto11$ on finite binary words. An infinite binary word either is a concatenation of
blocks $0$ and $11$, or it has a first occurrence of $10$ at a block boundary. So Cantor space is the disjoint union of the set $\mathcal C$ of
block sequences and the cylinders $Z_w=\mathsf c(w)10\{0,1\}^{\mathbb N}$, one for each finite word $w$. Index $Z_w$ by the dyadic
$x_w=0.w1\in\mathbb D$ in binary; this is a bijection. Let $f\in F$ have a tree-pair diagram with leaves $w_1,\dots,w_n$ and $w'_1,\dots,w'_n$,
so that $f(0.w_kz)=0.w'_kz$. Define $\hat f$ by the prefix replacements $\mathsf c(w_k)\mapsto\mathsf c(w'_k)$ and
$\mathsf c(v)10\mapsto\mathsf c(v')10$, where the interior vertices $v$ of the first tree are matched in order with the interior vertices $v'$ of
the second. Both families of prefixes are complete prefix codes, because splitting a leaf $w$ replaces $\mathsf c(w)$ by $\mathsf c(w0)$,
$\mathsf c(w)10$ and $\mathsf c(w1)$; expanding a matched pair of leaves adds a matched pair of gaps and changes nothing else, so $\hat f$ depends
only on $f$. It maps $Z_w$ onto $Z_{w'}$ with $x_{w'}=f(x_w)$: below a leaf this is the formula for $f$, and interior vertices correspond to the
breakpoints between leaves. On $\mathcal C$ it acts as $f$ acts on binary sequences, so $f\mapsto\hat f$ is an injective homomorphism $F\to V$.
Inside each $Z_w$ let $S_3$ permute the three cylinders $\mathsf c(w)10y\{0,1\}^{\mathbb N}$, $y\in\{0,10,11\}$. Copies with different $w$ have
disjoint supports, and $\hat f$ carries the copy at $x$ to the copy at $f(x)$. An element of the group they generate that acts trivially acts
trivially on $\mathcal C$, so its $F$-part is trivial, and then each of its $S_3$ coordinates is trivial.
\end{proof}

Let $a,b$ generate the copy of $S_3$ at $\frac12$, let $t\in F$ be
\[
t(x)=\begin{cases}2x,&0\le x\le\frac14,\\ x+\frac14,&\frac14\le x\le\frac12,\\ \frac{x+1}2,&\frac12\le x\le1,\end{cases}
\]
so that $t(\frac12)=\frac34$, and for $m\ge1$ let $\mathbb D_m$ be the set of the $2^m-1$ points of $\mathbb D$ with denominator at most
$2^m$.

\begin{proposition}\label{prop:thompson}
The group $\mathcal L$ is finitely presented, contains $F$, and contains no Houghton group $H_m$, so the bounds below do not come from the
Houghton subgroups of $V$. There are a finite set $R$ of relations of $\mathcal L$, containing $a^2,b^2,(ab)^3$, and a constant $C$ with the
following property: for every $m\ge1$ there are words $w_p$, $p\in\mathbb D_m$, with $w_p(\frac12)=p$ and
\[
\mathrm{Area}_R\bigl([w_pxw_p^{-1},w_qyw_q^{-1}]\bigr)\le Cm^2\qquad(p\ne q\in\mathbb D_m,\ x,y\in\{a,b\}).
\]
Consequently a chunk $E$ of $\mathcal L$, and hence of $V$, has $\log\sigma_E(r)\ge c\,r/\log^2r$ for all large $r$, and for $0<\Delta\le1$
every unitary assignment as in Theorem~\ref{thm:criterion} has $\log D\ge c_\Delta\,e^{-2}/\log^4(1/e)$ for all small $e$, where $c>0$ is a
constant and $c_\Delta>0$ depends only on $\Delta$.
\end{proposition}

\begin{proof}
\emph{Finite presentation.} Thompson's group $F$ is finitely presented, and it acts transitively on increasing pairs of points of
$\mathbb D$~\cite[Lemma~4.2]{CFP1996}, so it has finitely many orbits on $\mathbb D\times\mathbb D$. The stabilizer $F_{1/2}$ of $\frac12$ is the
product of the elements supported in $[0,\frac12]$ and in $[\frac12,1]$, two copies of $F$, so it is finitely generated. By Cornulier's
criterion for permutational wreath products, $\mathcal L$ is finitely presented; this action of $F$ is~\cite[Example~3.4]{Cornulier2006}. The bounds below do not use this.

\emph{Transport.} A standard dyadic interval is one of the form $[k2^{-j},(k+1)2^{-j}]$. If two partitions of $[0,1]$ into standard dyadic
intervals have the same number $n$ of intervals, the map sending the $i$th interval of the first affinely onto the $i$th interval of the second
lies in $F$ and has a tree-pair diagram with $n-1$ carets, and every standard dyadic interval can be split into any given number of standard
dyadic intervals. For $p<q$ in $\mathbb D_m$, the binary digits of $p$ split $[0,p]$ and $[p,1]$ into at most $m$ standard dyadic
intervals each, and $[p,q]$, cut at its point $c$ with the shortest binary expansion, splits into at most $m$ standard dyadic intervals
on each side of $c$. Hence there are $w_p\in F$ with $w_p(\frac12)=p$ and fewer than $2m$ carets, and $g_{pq}\in F$ with $g_{pq}(\frac12)=p$, $g_{pq}(\frac34)=q$ and fewer than $4m$
carets. The word length in $F$ is comparable to the number of carets of the reduced tree-pair diagram~\cite{BCS2001}, so these elements
are represented by words of length $O(m)$ in the generating set of the finite presentation of $F$ used in $R$ below; the Dehn functions of
two finite presentations of $F$ are equivalent, so Guba's quadratic bound holds for it. We write $w_p$ also for such a word. An element fixing
$\frac12$ maps $[0,\frac12]$ to itself, so the left and right subtrees of its reduced diagram are diagrams for its two restrictions, which are
elements of $F$ after affine rescaling, with fewer carets in total. Hence $F_{1/2}\cong F\times F$ is undistorted in $F$: its elements of
length $n$ in $F$ are words of length $O(n)$ in a fixed finite generating set $h_1,\dots,h_q$ of $F_{1/2}$.

\emph{Relations.} Let $R$ consist of a finite presentation of $F$ on a generating set that includes $t$ and the $h_j$, the words
$a^2,b^2,(ab)^3$, the commutators $[x,h_j]$, and the commutators $[x,tyt^{-1}]$, for $x,y\in\{a,b\}$. They hold in $\mathcal L$: each $h_j$
fixes $\frac12$, and the copies at $\frac12$ and at $t(\frac12)=\frac34$ commute.

\emph{Area.} Let $p<q$ in $\mathbb D_m$ and $g=g_{pq}$. The elements $g^{-1}w_p$ and $t^{-1}g^{-1}w_q$ fix $\frac12$ and have length
$O(m)$, so they are represented by words $\omega,\omega'$ of length $O(m)$ in the $h_j$. The words $w_p^{-1}g\omega$ and $w_q^{-1}gt\omega'$ are null
in $F$
and have length $O(m)$, so Guba's quadratic Dehn function~\cite{Guba2006} fills each with $O(m^2)$ cells. Replacing the four
occurrences of $w_p^{\pm1}$ in $[w_pxw_p^{-1},w_qyw_q^{-1}]$ by $(g\omega)^{\pm1}$, and the four of $w_q^{\pm1}$ by $(gt\omega')^{\pm1}$, therefore
costs
$O(m^2)$ cells. Moving $x^{\pm1}$ across $\omega$ and $y^{\pm1}$ across $\omega'$ with the relations $[x,h_j]$ costs $O(m)$ cells and leaves
$g[x,tyt^{-1}]g^{-1}$, a conjugate of a relation. For $p>q$ the commutator is the inverse of the one with the two copies exchanged. So every
area is $O(m^2)$.

\emph{Profile.} Apply Theorem~\ref{thm:criterion} with $N=2^m-1$ and $\Lambda=Cm^2$, where $C\ge1$, and let $l$ be the largest
length of a word of $R$. Given $r$, let $m$ be the largest
integer with $2^mm^2\le r/(10^8C(2l-1))$; then $2^m$ is at least of order $r/\log^2r$, and $\log\sigma_E(r)\ge(2^m-1)/16$. Given
$e$, let $m$ be the largest integer with $2^{m/2}m^2\le1/(CK_\Delta e)$; then $2^m$ is at least of order
$(CK_\Delta)^{-2}e^{-2}/\log^4(1/e)$, and $\log D\ge(2^m-1)\Delta^2/12$.

\emph{No Houghton subgroup.} A torsion element of $\mathcal L$ maps to a torsion element of the torsion-free group $F$, so it lies in
$\bigoplus_{\mathbb D}S_3$, which has exponent six. Every $H_m$ contains a five-cycle. Finally, $\mathcal L$ embeds in $V$ by
Lemma~\ref{lem:gap-coding}.
\end{proof}

The subgroup $\mathcal L$ contains $F$, whose soficity is open, so we do not know whether this chunk has permutation models at all; if it
does, they are very large. In particular, either $V$ is not sofic, or its models of this chunk need $2^{c\,r/\log^2r}$ points, close to the
largest bound that Theorem~\ref{thm:criterion} can give. If $F$ is amenable, so is $\mathcal L$, and then this chunk has finite profile. If $F$ is
hyperlinear, so is $\mathcal L$ (Proposition~\ref{prop:thompson-hyperlinear}), and then unitary models of the chunk exist at every defect. For
unitary models of Slofstra's group, his bound~\cite{Slofstra2018} gives $\log D\gtrsim e^{-c}$ for every $c<\frac12$, and, as Slofstra
notes~\cite[Corollary~1.3 and its proof]{Slofstra2018}, it gives permutation models with $\log\sigma\gtrsim r^c$ for every $c<\frac14$; the group $\mathcal L$ gives $e^{-2}$ and $r$ up to logarithms. The
group $\mathcal L$ has the form $\bigl(\bigoplus_{\Gamma/\Gamma_0}A\bigr)\rtimes\Gamma$ of the non-sofic wreath products of Kun and Thom~\cite{KunThom2026},
which are built from groups with property (T); here the input is instead cheap transport of non-abelian copies of $S_3$.

Soficity and hyperlinearity pass to wreath products whose acting group is sofic~\cite{HayesSale2018}, and to permutational wreath products
whose acting group is sofic (respectively hyperlinear) and whose action is sofic~\cite[Theorems~3.6 and~3.8]{GKEP2025},~\cite{AlekseevBradford2026}.
Since $F$ acts on $\mathbb D$ transitively and faithfully, a sofic action of $F$ on $\mathbb D$ would already make $F$
sofic~\cite[Theorem~2.12]{GKEP2025}, so these results give nothing here; the next proposition assumes only that $F$ is hyperlinear.

\begin{proposition}\label{prop:thompson-hyperlinear}
The group $\mathcal L$ is hyperlinear if and only if $F$ is.
\end{proposition}

\begin{proof}
Subgroups of hyperlinear groups are hyperlinear, so suppose that $F$ is hyperlinear. A countable group $\Gamma$ is hyperlinear if and only if
its von Neumann algebra $L(\Gamma)$ has a trace-preserving embedding into an ultrapower $\mathcal R^\omega$ of the hyperfinite
$\mathrm{II}_1$ factor~\cite{CapraroLupini2015}, and then every separable von Neumann subalgebra of $L(\Gamma)^\omega$ embeds in
$\mathcal R^\omega$ as well, by a diagonal argument. It therefore suffices to embed $L(\mathcal L)$ into $L(F)^\omega$ preserving the trace,
for a free ultrafilter $\omega$ on $\mathbb N$. We write $\tau$ for the canonical traces and $F_J$ for the subgroup of elements supported in an
interval $J$. If $J,J'$ are disjoint, $L(F_J)$ and $L(F_{J'})$ commute and $\tau(xy)=\tau(x)\tau(y)$ for $x\in L(F_J)$ and $y\in L(F_{J'})$,
since $F_J$ and $F_{J'}$ generate their direct product. If $f,f'\in F$ agree on $J$, then $\lambda(f)x\lambda(f)^{-1}=\lambda(f')x\lambda(f')^{-1}$
for $x\in L(F_J)$, because $f'^{-1}f$ fixes $J$ pointwise and therefore commutes with $F_J$.

\emph{A localized copy of $S_3$.} Choose $\sigma\in F$ supported in $[\frac12,\frac34]$, with $\sigma(x)<x$ on $(\frac12,\frac34)$ and
$\sigma(x)=\frac12+\frac12(x-\frac12)$ near $\frac12$, and a dyadic $x_0\in(\frac12,\frac34)$. The intervals
$J_k=\sigma^k((\sigma(x_0),x_0))$, $k\in\mathbb Z$, are disjoint and accumulate at $\frac12$ as $k\to\infty$. For a nonempty open interval $J$ with
dyadic endpoints, $F_J\cong F$ has infinite conjugacy classes, since a nontrivial element has conjugates with infinitely many different
supports, so $L(F_J)$ is a $\mathrm{II}_1$ factor. In $L(F_{J_0})$ choose a unital copy of $M_6$ and, in its relative commutant, which is again a
$\mathrm{II}_1$ factor, a self-adjoint $Z_0$ with uniform distribution on $[0,1]$. Let $v_0:S_3\to U(M_6)$ be the regular representation, so
that $\tau(v_0(g))=0$ for $g\ne1$, and put $v_k=\lambda(\sigma)^kv_0\lambda(\sigma)^{-k}$ and $Z_k=\lambda(\sigma)^kZ_0\lambda(\sigma)^{-k}$,
which lie in $L(F_{J_k})$. By the factorization of traces, the $Z_k$ are independent and uniform. For a finite $K\subset\mathbb Z$ and $k\in K$,
let $p^K_k$ be the spectral projection of the event that $Z_k$ is the largest of the $Z_j$, $j\in K$. Ties have trace zero, so these projections
are orthogonal and sum to $1$, and they commute with every $v_j(g)$. Hence
\[
V_K(g)=\sum_{k\in K}p^K_kv_k(g)
\]
is a representation of $S_3$. Since $v_k(g)$ lies in a copy of $M_6$ whose relative commutant contains $Z_k$, and the other $Z_j$ lie in
algebras with factorized trace, $\tau(p^K_kv_k(g))=\tau(p^K_k)\tau(v_0(g))$; so $\tau(V_K(g))=0$ for $g\ne1$. If $K$ is an interval of $n$ integers and $|q|\le n$,
the two selected indices agree unless the largest $Z_j$ over $K\cup(K+q)$ has its index outside $K\cap(K+q)$, an event of trace
$2|q|/(n+|q|)$, and on its complement $V_K(g)=V_{K+q}(g)$. Hence $\|V_{K+q}(g)-V_K(g)\|_2^2\le8|q|/(n+|q|)$.

\emph{Invariance.} Put $V_n=V_{\{n,\dots,2n-1\}}$, and let $h\in F_{1/2}$. Near $\frac12$ on the right, $h(x)=\frac12+2^m(x-\frac12)$ for an
integer $m$, so $h$ agrees with $\sigma^{-m}$ on $J_k$ for all large $k$. Conjugation by $\lambda(h)$ therefore carries $v_k$ and $Z_k$ to
$v_{k-m}$ and $Z_{k-m}$ for large $k$, and $V_n$ to $V_{\{n,\dots,2n-1\}-m}$ for large $n$, so $\|\lambda(h)V_n(g)\lambda(h)^{-1}-V_n(g)\|_2\to0$.
Thus $V(g)=(V_n(g))_\omega$ is a representation of $S_3$ in $L(F)^\omega$ with $\tau(V(g))=0$ for $g\ne1$, and it commutes with $\lambda(F_{1/2})$.

\emph{The embedding.} For $x\in\mathbb D$ choose $f_x\in F$ with $f_x(\frac12)=x$ and put $V_x=\lambda(f_x)V\lambda(f_x)^{-1}$. By the
invariance this does not depend on the choice of $f_x$, and $\lambda(f)V_x\lambda(f)^{-1}=V_{f(x)}$ for $f\in F$. The representatives $V_n(g)$ lie
in $L(F_{I_n})$, where $I_n=\bigcup_{n\le k<2n}J_k$ shrinks to $\frac12$ from the right, so those of $V_x$ lie in $L(F_{f_x(I_n)})$, and for $x\ne y$
the intervals $f_x(I_n)$ and $f_y(I_n)$ are disjoint for large $n$. So copies at distinct points commute, and the $V_x$ together with $\lambda$
define a homomorphism $\pi:\mathcal L\to U(L(F)^\omega)$. Let $\gamma=\bigl(\prod_{x\in X}\ell_x\bigr)f$, with $X$ finite, $\ell_x$ in the copy of $S_3$
at $x$ and $f\in F$. If $f=1$ and some $\ell_x\ne1$, the factors lie in commuting algebras with factorized trace for large $n$, so
$\tau(\pi(\gamma))=\prod_x\tau(V(\ell_x))=0$. If $f\ne1$, choose a point moved by $f$ outside $X$. For large $n$ it lies outside the closure of
every $f_x(I_n)$, so $f$ is not in the group $K_n$ generated by the $F_{f_x(I_n)}$, and $\tau(a\lambda(f))=0$ for every $a\in L(K_n)$, in
particular for the representatives of $\prod_x\pi(\ell_x)$. Hence $\tau(\pi(\gamma))=0$ whenever $\gamma\ne1$, and $\pi$ extends to a
trace-preserving embedding of $L(\mathcal L)$ into $L(F)^\omega$.
\end{proof}

\subsection{From orbits to configurations}\label{sec:toward}

The two steps isolate what is needed: exponentially many sites reachable by short words, pairwise commutators of polylogarithmic area, and an
amenable ambient group. In $\mathcal G$ the orbit of $0$ grows exponentially, but we do not know how to commute two copies at points of
opposite parity cheaply. In $\mathcal L$ everything works except amenability, which is the open question of whether $F$ is amenable. The
configuration lamps of the next section obtain all three at once by a change of viewpoint: the sites are finite binary configurations of the points of Houghton's three rays
rather than points, the transport is by affine maps that touch a bounded number of coordinates at a time, and the ambient group is locally
finite-by-$\Z^2$.
